После восстановления виртуальной машины из снимка её Git-репозиторий может существенно отставать от актуального состояния на удалённых серверах. Это происходит, потому что «откат во времени» возвращает и состояние всех файлов, которое было в момент создания снапшота. Чтобы синхронизировать код и продолжить работу, нужно выполнить несколько шагов по получению всех последних изменений.

В начале следует ввести команду \texttt{git fetch --all}. Она связывается со всеми настроенными удалёнными репозиториями и загружает информацию о новых коммитах, ветках и тегах. Важно, что эта команда не изменяет локальные файлы и текущую ветку. Это позволяет сначала просто посмотреть, какие изменения произошли, и спланировать их интеграцию.

В проекте «Час ЕГЭ» основной веткой для разработки является devel. Чтобы получить и применить изменения из официального репозитория, обычно используется команда git pull origin devel. Она автоматически выполняет \texttt{git fetch origin devel}, а затем пытается слить (merge) загруженные изменения с локальной веткой devel.

Репозиторий myfork — это наша собственная копия на GitHub. Команда \texttt{git pull myfork devel} используется, чтобы получить новые коммиты, которые мы или другие могли записать в наш форк.

После того как мы подтянули все новые коммиты, запуск \texttt{grunt} пересоберёт проект заново, применив все полученные изменения и обновив рабочие файлы.